\section{Generating Sample Paths}

\subsection{Brownian Motion}

\subsubsection*{One-dimensional BM}

A standard Brownian motion $W(t)$ on $[0,T]$ satisfies $W(0)=0$, continuous sample paths, independent increments, and
$W(t)-W(s)\sim N(0,t-s), 0\le s<t\le T$.

Hence $W(t)\sim N(0,t)$. A Brownian motion with drift $\mu$ and diffusion coefficient $\sigma^2$, written $X\sim\operatorname{BM}(\mu,\sigma^2)$, can be represented by
$X(t)=X(0)+\mu t+\sigma W(t)$,

so $X(t)\sim N(X(0)+\mu t,\sigma^2t)$ and
$dX(t)=\mu\,dt+\sigma\,dW(t)$.

With deterministic time-varying coefficients,

$dX(t)=\mu(t)\,dt+\sigma(t)\,dW(t)$,
$X(t)=X(0)+\int_0^t\mu(s)\,ds+\int_0^t\sigma(s)\,dW(s)$.

For $0\le s<t$,
$\mathbb{E}[X(t)-X(s)]=\int_s^t\mu(u)\,du$,
$\operatorname{Var}[X(t)-X(s)]=\int_s^t\sigma^2(u)\,du$.

\subsubsection*{Random-walk construction}

For dates $0=t_0<t_1<\cdots<t_n$, let $Z_1,\ldots,Z_n$ be independent $N(0,1)$. For standard BM,
$W(t_{i+1})=W(t_i)+\sqrt{t_{i+1}-t_i}\,Z_{i+1},\qquad i=0,\ldots,n-1$.

For $X\sim\operatorname{BM}(\mu,\sigma^2)$ with constant coefficients,

$X(t_{i+1})=X(t_i)+\mu(t_{i+1}-t_i)+\sigma\sqrt{t_{i+1}-t_i}\,Z_{i+1}$.

For time-dependent coefficients, the exact grid recursion is

$X(t_{i+1})=X(t_i)+\int_{t_i}^{t_{i+1}}\mu(s)\,ds+\sqrt{\int_{t_i}^{t_{i+1}}\sigma^2(u)\,du}\,Z_{i+1}$.

These recursions are exact at the grid points. Discretization error appears if the path is deterministically interpolated between grid points. The Euler approximation
$X(t_{i+1})=X(t_i)+\mu(t_i)(t_{i+1}-t_i)+\sigma(t_i)\sqrt{t_{i+1}-t_i}\,Z_{i+1}$

is generally biased even at the grid dates when $\mu,\sigma$ vary in time, since it uses approximate increment mean and variance.

\subsubsection*{Multivariate normal view}

The vector $(W(t_1),\ldots,W(t_n))$ is multivariate normal with mean $0$ and covariance matrix $C$ satisfying
$C_{ij}=\operatorname{Cov}(W(t_i),W(t_j))=\min(t_i,t_j)$.

Thus one may simulate it as $AZ$, where $Z\sim N(0,I_n)$ and $AA^\top=C$. For the Brownian covariance matrix, a lower triangular factor is compactly

$A_{ij}=\sqrt{t_j-t_{j-1}}\,\mathbf{1}_{\{i\ge j\}},\quad t_0=0$.
For $X\sim\operatorname{BM}(\mu,\sigma^2)$, the mean vector is $\mu t_i$, the covariance matrix is $\sigma^2C$, and the Cholesky factor is $\sigma A$.

\subsection{Multidimensional Brownian Motion}

\subsubsection*{Definition and representation}

A standard $d$-dimensional Brownian motion $W(t)=(W_1(t),\ldots,W_d(t))^\top$ has

$W(t)-W(s)\sim N(0,(t-s)I)$.

Thus each coordinate is standard BM and different coordinates are independent. For $\mu\in\mathbb{R}^d$ and positive semidefinite $\Sigma$, $X\sim\operatorname{BM}(\mu,\Sigma)$ means

$X(t)-X(s)\sim N((t-s)\mu,(t-s)\Sigma)$.

If $BB^\top=\Sigma$ and $W$ is standard BM, then
$X(t)=X(0)+\mu t+BW(t)$,

and $dX(t)=\mu\,dt+B\,dW(t)$. The law depends on $B$ only through $BB^\top$.

\subsubsection*{Time-varying coefficients}

For
$dX(t)=\mu(t)\,dt+B(t)\,dW(t),\qquad B(t)B(t)^\top=\Sigma(t)$,

the increment distribution is
$X(t)-X(s)\sim N\left(\int_s^t\mu(u)\,du,\int_s^t\Sigma(u)\,du\right)$.

For $X\sim\operatorname{BM}(\mu,\Sigma)$,
$\operatorname{Cov}[X_i(s),X_j(t)]=\min(s,t)\Sigma_{ij}$.

\subsubsection*{Simulation}

Let $Z_1,Z_2,\ldots$ be independent $N(0,I)$ vectors. Standard $d$-dimensional BM is simulated by
$W(t_{i+1})=W(t_i)+\sqrt{t_{i+1}-t_i}\,Z_{i+1}$.

For $X\sim\operatorname{BM}(\mu,\Sigma)$, factor $\Sigma=BB^\top$ and use

$X(t_{i+1})=X(t_i)+\mu(t_{i+1}-t_i)+\sqrt{t_{i+1}-t_i}\,BZ_{i+1}$.

For time-dependent covariance, use
$X(t_{i+1})=X(t_i)+\int_{t_i}^{t_{i+1}}\mu(s)\,ds+B_iZ_{i+1}$,

where
$B_iB_i^\top=\int_{t_i}^{t_{i+1}}\Sigma(u)\,du$.

\subsection{Geometric Brownian Motion}

\subsubsection*{Exact solution and recursion}

A geometric Brownian motion is exponential Brownian motion:

$\frac{dS(t)}{S(t)}=\mu\,dt+\sigma\,dW(t),\qquad S\sim\operatorname{GBM}(\mu,\sigma^2)$.

The exact solution is
$S(t)=S(0)\exp\left(\left[\mu-\frac12\sigma^2\right]t+\sigma W(t)\right)$.

For $u<t$,
$S(t)=S(u)\exp\left(\left[\mu-\frac12\sigma^2\right](t-u)+\sigma(W(t)-W(u))\right)$.

Hence for $0=t_0<t_1<\cdots<t_n$,

$S(t_{i+1})=S(t_i)\exp\left(\left[\mu-\frac12\sigma^2\right]\Delta t_i+\sigma\sqrt{\Delta t_i}\,Z_{i+1}\right)$,
where $\Delta t_i=t_{i+1}-t_i$. 
\subsubsection*{Lognormal distribution}

Write $Y\sim LN(\mu,\sigma^2)$ if $Y=\exp(\mu+\sigma Z)$ with $Z\sim N(0,1)$. Then

$P(Y\le y)=\Phi\left(\frac{\log y-\mu}{\sigma}\right)$,
$f(y)=\frac{1}{y\sigma\sqrt{2\pi}}\exp\left(-\frac{(\log y-\mu)^2}{2\sigma^2}\right)$.

Using $\mathbb{E}[e^{aZ}]=e^{a^2/2}$,
$\mathbb{E}[Y]=e^{\mu+\sigma^2/2},\qquad \operatorname{var}(Y)=e^{2\mu+\sigma^2}(e^{\sigma^2}-1)$.

Also $P(Y\le e^\mu)=1/2$, so $e^\mu$ is the median. For $S\sim\operatorname{GBM}(\mu,\sigma^2)$,

$\frac{S(t)}{S(0)}\sim LN\left(\left[\mu-\frac12\sigma^2\right]t,\sigma^2t\right)$,

$\mathbb{E}[S(t)]=e^{\mu t}S(0),\quad \operatorname{var}[S(t)]=e^{2\mu t}S^2(0)(e^{\sigma^2t}-1)$.

The Markov conditional mean is

$\mathbb{E}[S(t)\mid S(\tau),0\le\tau\le u]=\mathbb{E}[S(t)\mid S(u)]=e^{\mu(t-u)}S(u)$.

\subsection{Risk-Neutral Dynamics}

\subsubsection*{Assets and futures}

With constant continuously compounded rate $r$, $\beta(t)=e^{rt}$. Under the risk-neutral measure, for a non-dividend-paying asset,

$\frac{S(u)}{\beta(u)}=\mathbb{E}\left[\frac{S(t)}{\beta(t)}\mid \{S(\tau),0\le\tau\le u\}\right]$.

Thus the risk-neutral GBM dynamics are
$\frac{dS(t)}{S(t)}=r\,dt+\sigma\,dW(t)$.

With continuous dividend yield $\delta$,
$\frac{dS(t)}{S(t)}=(r-\delta)\,dt+\sigma\,dW(t)$.

For a futures price $F(t,T)$ with fixed maturity $T$, the risk-neutral futures price is a martingale. If modeled by GBM,
$\frac{dF(t,T)}{F(t,T)}=\sigma\,dW(t)$.

\subsection{Path-Dependent Payoffs}

\subsubsection*{Asian options}

For monitoring dates $t_1,\ldots,t_n$, the arithmetic average is
$\bar{S}=\frac{1}{n}\sum_{i=1}^n S(t_i)$.

Asian call and put payoffs are $(\bar{S}-K)^+$ and $(K-\bar{S})^+$. Floating-strike variants include $(S(T)-\bar{S})^+$ and $(\bar{S}-S(T))^+$. The arithmetic average distribution is generally intractable, so simulation is natural.

For continuous monitoring over $[u,t]$,
$\bar{S}=\frac{1}{t-u}\int_u^tS(\tau)\,d\tau$.

% \subsubsection*{Geometric average options}

% The discrete geometric average is
% $G=\left(\prod_{i=1}^nS(t_i)\right)^{1/n}$.

% Because lognormal variables remain lognormal under geometric averaging, these options are analytically convenient. Under risk-neutral GBM with dividend yield $\delta$,
% $G=S(0)\exp\left(\left[r-\delta-\frac12\sigma^2\right]\frac{1}{n}\sum_{i=1}^nt_i+\frac{\sigma}{n}\sum_{i=1}^nW(t_i)\right)$.

% Since $\operatorname{Cov}(W(s),W(t))=\min(s,t)$,
% $\sum_{i=1}^nW(t_i)\sim N\left(0,\sum_{i=1}^n(2i-1)t_{n+1-i}\right)$.

% The geometric average has the same distribution as a GBM value at

% $T_G=\frac{1}{n}\sum_{i=1}^nt_i$,
% with adjusted volatility and dividend yield

% $\bar{\sigma}^2=\frac{\sigma^2}{n^2T_G}\sum_{i=1}^n(2i-1)t_{n+1-i}$,
% $\bar{\delta}=\delta+\frac12\sigma^2-\frac12\bar{\sigma}^2$.

% Thus a geometric-average option can be priced using Black-Scholes with adjusted parameters. The continuous geometric average

% $\exp\left(\frac{1}{t-u}\int_u^t\log S(\tau)\,d\tau\right)$
% is also lognormally distributed.

% \subsubsection*{Barrier and lookback options}

% A discretely monitored down-and-out call with barrier $b$ has payoff

% $\mathbf{1}_{\{\tau(b)>T\}}(S(T)-K)^+,\qquad \tau(b)=\inf\{t_i:S(t_i)<b\}$.

% If no grid value crosses below $b$, set $\tau(b)=\infty$. A down-and-in call has payoff

% $\mathbf{1}_{\{\tau(b)\le T\}}(S(T)-K)^+$.

% For continuous monitoring, use $\tilde{\tau}(b)=\inf\{t\ge0:S(t)\le b\}$. Discretely monitored barrier options are priced by simulating $S(t_1),\ldots,S(t_n),S(T)$.

% Lookback options depend on path extrema. Discrete lookback put and call payoffs at $t_n$ are
% $\max_{i=1,\ldots,n}S(t_i)-S(t_n),\qquad S(t_n)-\min_{i=1,\ldots,n}S(t_i)$.

\subsection{Multi-Asset GBM}

\subsubsection*{Correlated assets}

Let $X_i$ be standard Brownian motions with correlations $\rho_{ij}$. A multidimensional GBM is
$\frac{dS_i(t)}{S_i(t)}=\mu_i\,dt+\sigma_i\,dX_i(t),\qquad i=1,\ldots,d$.

Define $\Sigma_{ij}=\sigma_i\sigma_j\rho_{ij}$. Then $(\sigma_1X_1,\ldots,\sigma_dX_d)\sim\operatorname{BM}(0,\Sigma)$ and $S=(S_1,\ldots,S_d)\sim\operatorname{GBM}(\mu,\Sigma)$.

The price covariance is

$\operatorname{Cov}[S_i(t),S_j(t)]=S_i(0)S_j(0)e^{(\mu_i+\mu_j)t}\left(e^{\rho_{ij}\sigma_i\sigma_jt}-1\right)$.

If $AA^\top=\Sigma$ and $a_i$ is the $i$th row of $A$,

$\frac{dS_i(t)}{S_i(t)}=\mu_i\,dt+a_i\,dW(t)=\mu_i\,dt+\sum_{j=1}^dA_{ij}\,dW_j(t)$.

\subsubsection*{Simulation and payoffs}

With independent $Z_k=(Z_{k1},\ldots,Z_{kd})\sim N(0,I)$,

$S_i(t_{k+1})=S_i(t_k)\exp\left(\left[\mu_i-\frac12\sigma_i^2\right]\Delta t_k+\sqrt{\Delta t_k}\sum_{j=1}^dA_{ij}Z_{k+1,j}\right)$.

Under the risk-neutral measure, $\mu_i=r-\delta_i$. 
